\documentclass[11pt,a4paper]{article}

\usepackage[a4paper, margin=1in]{geometry}
\usepackage{amsmath, amssymb}
\usepackage{graphicx}
\usepackage{caption}
\usepackage{subcaption}
\usepackage{float}
\usepackage{hyperref}
\usepackage{booktabs}
\usepackage{array}
\usepackage{longtable}

\hypersetup{
    colorlinks=true,
    linkcolor=blue,
    urlcolor=blue
}

\title{
    IC253: Data Structures and Algorithm \\
    Assignment 4 Report
}
\author{Name: Daksh Rathi \\ Roll No: B25132}
\date{}

\begin{document}

\maketitle

\thispagestyle{empty}

% =====================================================================
\section{Introduction}
% =====================================================================

This report presents the implementation of both tasks of Assignment 4.
Task 1 implements a Smart City Road Network Optimizer on an undirected
weighted graph with two edge weights (construction cost and travel time),
and Task 2 implements a Logistics Dispatch Centre that combines a live
max-heap priority queue with an in-place randomised quicksort for shift-end
record sorting. The report summarises the validation policies, the main
data structures, the implementation approach, time--space complexity, and
the test-case outputs for both tasks.

% =====================================================================
\section{Task 1: Smart City Road Network Optimizer}
% =====================================================================

\subsection{Objective}
The objective of this task is to model a city as an undirected weighted
graph with \(N\) intersections and \(M\) roads, where every road carries a
construction cost and a travel time. The class \texttt{CityGraph} supports
graph validation, MST construction using Kruskal's algorithm, shortest
travel-time path computation using Dijkstra's algorithm, comparison of the
MST path against the optimal shortest path, and identification of the
critical roads of the MST.

\subsection{Approach}
The complete graph is managed by the class \texttt{CityGraph}. Each road is
stored in an \texttt{Edge} struct containing \texttt{u}, \texttt{v},
\texttt{cost}, and \texttt{travelTime}. The input edge list is kept in
\texttt{edges} and is never modified after the constructor. The full-graph
adjacency chain \texttt{adjFull} is built once inside the constructor using
\texttt{buildFullAdjacency()}, which means that \texttt{shortestPath()} can
be called on its own without requiring \texttt{buildMST()} to have been
invoked first. The MST-only adjacency chain \texttt{adjMST} is built later
inside \texttt{buildMST()} via \texttt{buildMSTAdjacency()}, and a boolean
flag \texttt{mstBuilt} is set to true at the end of \texttt{buildMST()};
\texttt{compareNetworks()} and \texttt{criticalRoads()} check this flag
before running.

The MST is constructed by Kruskal's algorithm inside \texttt{buildMST()}.
The edge indices are sorted in ascending order of \texttt{cost} using
\texttt{sortEdgeIndicesByCost()}, which is a plain selection sort on
indices so the original input order is preserved. A Disjoint Set Union
(DSU) with path compression and union by size is maintained in the arrays
\texttt{parent} and \texttt{setSize} through the helpers \texttt{initDSU()},
\texttt{findParent()}, and \texttt{unionSets()}. Each edge is either
accepted (when its endpoints lie in different DSU components) or rejected
(when they already lie in the same component and the edge would form a
cycle), and the running MST cost is printed after every acceptance. The
same DSU is reused inside \texttt{criticalRoads()}, which reruns Kruskal's
pass once per MST edge while skipping that edge to check whether the MST
cost increases or the graph becomes disconnected.

Shortest paths are computed by the helper \texttt{dijkstra()} in the simple
\(O(N^{2})\) form: at every step the unvisited node with the smallest
tentative distance is picked by a linear scan of the \texttt{dist} array,
its outgoing edges are relaxed, and the predecessor is recorded in the
\texttt{prev} array. The actual path from \(S\) to \(T\) is then rebuilt by
\texttt{reconstructPath()}. The helper \texttt{getPathTravelTime()} is a
thin wrapper that returns just the integer travel time, and it is used
twice inside \texttt{compareNetworks()} --- once on \texttt{adjFull} and
once on \texttt{adjMST}. Connectivity during input validation is checked by
\texttt{countReachableFromZero()}, which runs a BFS from node 0 using a
\texttt{vector<int>} with a front index in place of a queue.

\subsection{Graph Validation Policy}
The validation policy enforced by \texttt{validateGraph()} is:
\begin{itemize}
    \item Self-loops are forbidden: an edge with \(u = v\) is rejected.
    \item Every node id used by an edge must lie in \([0, N - 1]\).
    \item Construction cost and travel time must both be non-negative.
    \item Duplicate roads between the same unordered pair of intersections are forbidden.
    \item The graph must be connected: a BFS from node 0 through \texttt{countReachableFromZero()} must reach every node.
\end{itemize}
The rules are checked in the order listed above, and the first rule that
fails rejects the input with a specific reason. Invalid source and
destination node ids inside \texttt{shortestPath()} and
\texttt{compareNetworks()} are caught at the start of the call at runtime.

\subsection{Functions and Implementation}
\texttt{initDSU(n)}, \texttt{findParent(x)}, and \texttt{unionSets(x, y)}
implement the DSU with path compression and union by size.
\texttt{sortEdgeIndicesByCost(indexList)} is a selection sort on the edge
indices. \texttt{dijkstra(adj, S, dist, prev)} computes single-source
shortest paths in \(O(N^{2})\). \texttt{reconstructPath(T, prev, dist)}
rebuilds the \(S \to T\) path from \texttt{prev}.
\texttt{getPathTravelTime(adj, S, T)} returns just the integer travel time.
\texttt{buildFullAdjacency()} fills \texttt{adjFull} once in the
constructor, and \texttt{buildMSTAdjacency()} fills \texttt{adjMST} at the
end of \texttt{buildMST()}. \texttt{countReachableFromZero()} performs the
BFS-based connectivity check. The public operations
\texttt{validateGraph()}, \texttt{buildMST()}, \texttt{shortestPath(S, T)},
\texttt{compareNetworks(S, T)}, and \texttt{criticalRoads()} orchestrate
the helpers as described above.

\subsection{Time and Space Complexity Analysis}
Let \(N\) be the number of intersections, \(M\) the number of roads, and
\(E = N - 1\) the number of MST edges of a valid input.

\begin{table}[H]
\centering
\begin{tabular}{|p{4.3cm}|p{3.3cm}|p{2.5cm}|p{5.2cm}|}
\hline
\textbf{Operation} & \textbf{Time Complexity} & \textbf{Space Complexity} & \textbf{Explanation} \\
\hline
\texttt{findParent()},
\texttt{unionSets()}
& nearly \(O(1)\) amortised & \(O(N)\) &
Path compression and union by size keep the DSU operations almost constant-time. \\
\hline
\texttt{sortEdgeIndicesByCost()}
& \(O(M^{2})\) & \(O(M)\) &
Selection sort on the \(M\) edge indices. \\
\hline
\texttt{dijkstra()}
& \(O(N^{2} + M)\) & \(O(N)\) &
At each iteration the next node is picked by a linear scan of \texttt{dist}. \\
\hline
\texttt{reconstructPath()}
& \(O(N)\) & \(O(N)\) &
Walks the \texttt{prev} array from \(T\) back to the source. \\
\hline
\texttt{countReachableFromZero()}
& \(O(N + M)\) & \(O(N + M)\) &
Single BFS over the unweighted version of the graph. \\
\hline
\texttt{validateGraph()}
& \(O(M^{2} + N + M)\) & \(O(N + M)\) &
The duplicate-edge double loop is the dominant \(O(M^{2})\) term. \\
\hline
\texttt{buildMST()}
& \(O(M^{2})\) & \(O(N + M)\) &
Dominated by the selection sort on edge indices. \\
\hline
\texttt{shortestPath(S, T)}
& \(O(N^{2} + M)\) & \(O(N)\) &
One \texttt{dijkstra()} call followed by path reconstruction. \\
\hline
\texttt{compareNetworks(S, T)}
& \(O(N^{2} + M)\) & \(O(N)\) &
Two \texttt{dijkstra()} calls, one on \texttt{adjFull} and one on \texttt{adjMST}. \\
\hline
\texttt{criticalRoads()}
& \(O(E \cdot M^{2})\) & \(O(N + M)\) &
Kruskal's pass is rerun once per MST edge. \\
\hline
\end{tabular}
\caption{Task 1 time and space complexity.}
\end{table}

Total storage used by the class is \(O(N + M)\) because of \texttt{edges},
\texttt{mstEdges}, \texttt{adjFull}, and \texttt{adjMST}.

\subsection{Test Case Demonstrations}

\begin{figure}[H]
\centering
\begin{subfigure}[b]{0.48\textwidth}
    \includegraphics[width=\textwidth,keepaspectratio]{task1_tc1.png}
    \caption{Task 1 -- Test Case 1 Output}
\end{subfigure}
\hfill
\begin{subfigure}[b]{0.48\textwidth}
    \includegraphics[width=\textwidth,keepaspectratio]{task1_tc2.png}
    \caption{Task 1 -- Test Case 2 Output}
\end{subfigure}

\vspace{0.6em}

\begin{subfigure}[b]{0.48\textwidth}
    \includegraphics[width=\textwidth,keepaspectratio]{task1_tc3.png}
    \caption{Task 1 -- Test Case 3 Output}
\end{subfigure}
\hfill
\begin{subfigure}[b]{0.48\textwidth}
    \includegraphics[width=\textwidth,keepaspectratio]{task1_tc4.png}
    \caption{Task 1 -- Test Case 4 Output}
\end{subfigure}
\caption{Task 1 test cases}
\end{figure}

% =====================================================================
\section{Task 2: Logistics Dispatch Centre}
% =====================================================================

\subsection{Objective}
The objective of this task is to implement a regional disaster-relief
dispatch centre that stores \(N\) aid requests, each with a unique
\texttt{requestID}, an integer \texttt{priority} in \([0, 100]\), and a
submission \texttt{timestamp}. The class \texttt{DispatchCentre} supports
input validation, live heap-based dispatch, in-place priority updates, a
top-\(k\) peek that does not modify the live structure, and shift-end
in-place sorting of records by either \texttt{priority} or
\texttt{timestamp}.

\subsection{Approach}
The complete dispatch centre is managed by the class
\texttt{DispatchCentre}. Each request is stored in a \texttt{Request}
struct containing \texttt{requestID}, \texttt{priority}, and
\texttt{timestamp}. The max-heap itself is kept inside a single data member
\texttt{heap} of type \texttt{vector<Request>}, treated as a one-dimensional
array with the standard indexing \(\text{parent} = (i - 1) / 2\),
\(\text{leftChild} = 2i + 1\), \(\text{rightChild} = 2i + 2\). The internal
helpers \texttt{swapNodes()}, \texttt{siftUp()}, and \texttt{siftDown()}
restore the max-heap invariant after any single change.

The heap is built inside \texttt{buildStructure()} using bottom-up heapify
(Floyd's algorithm): starting from the last non-leaf node at index
\(n / 2 - 1\) and walking backwards to the root, each subtree is heapified
by a sift-down. This gives overall \(O(n)\) construction instead of the
\(O(n \log n)\) of repeated insertion. \texttt{dispatchNext()} removes the
root, replaces it with the last element, shrinks the heap by one, and sifts
the new root down. \texttt{updatePriority()} locates the target request
with \texttt{findIndexByID()}, which is a linear scan over the heap array,
then overwrites the priority and sifts the node either up or down depending
on whether the new priority is greater or smaller than the old one.

Shift-end sorting is handled by \texttt{sortRecords()}, which calls the
recursive helper \texttt{quicksortHelper()}. This is a randomised Lomuto
quicksort: at each recursive call a pivot index is chosen uniformly at
random in \([\text{lo}, \text{hi}]\), the pivot is moved to the right end,
and the array is partitioned by pushing elements greater than or equal to
the pivot to the left. The sort is entirely in-place --- the only extra
memory is the recursion stack, \(O(\log n)\) on average. The comparison
returns ``greater or equal'' rather than ``less or equal'', so the final
order is \emph{descending} on the chosen key. This behaviour is stated
explicitly in the banner printed by \texttt{sortRecords()}. The
\texttt{sortRecords()} function supports both \texttt{"priority"} and
\texttt{"timestamp"} via the helper \texttt{getKey()}.

The top-\(k\) peek \texttt{topKRequests(k)} is implemented on a local copy
\texttt{tempHeap} so the live \texttt{heap} is never touched. The function
performs \(k\) extract-max operations on \texttt{tempHeap}, prints each
extracted request, and \emph{returns} the collected list as a
\texttt{vector<Request>} in decreasing order of priority.

\subsection{Data Handling Policy}
The data handling policy enforced by \texttt{validateInput()} and
\texttt{updatePriority()} is:
\begin{itemize}
    \item Every \texttt{requestID} inside one input batch must be unique.
    \item \texttt{priority} must be an integer in \([0, 100]\).
    \item Duplicate request IDs are rejected at input time, with the offending ID reported.
    \item Out-of-range priorities are rejected at input time, with the offending value reported.
    \item Inside \texttt{updatePriority()}, an unknown \texttt{requestID} and an out-of-range new priority are both rejected with a clear error, and the heap is left unchanged in either case.
\end{itemize}

\subsection{Functions and Implementation}
\texttt{swapNodes(i, j)}, \texttt{siftUp(i)}, and \texttt{siftDown(i, n)}
restore the max-heap invariant after any single change.
\texttt{findIndexByID(reqID)} is a linear scan that returns the heap index
of a request given its ID, or \(-1\) if the ID is absent.
\texttt{getKey(r, key)} returns either \texttt{r.priority} or
\texttt{r.timestamp} depending on the \texttt{key} string and is used by
the partition step. \texttt{quicksortHelper(arr, lo, hi, key, partitionCount)}
implements the recursive randomised Lomuto quicksort.
\texttt{printHeap()} prints the live heap as an array after every
modification.

The public operations are \texttt{validateInput(requests)},
\texttt{buildStructure(requests)}, \texttt{dispatchNext()},
\texttt{updatePriority(reqID, newPriority)}, \texttt{sortRecords(records, key)},
and \texttt{topKRequests(k)}. The last one returns a
\texttt{vector<Request>} containing the top \(k\) requests in decreasing
order of priority; the same list is also printed.

\subsection{Time and Space Complexity Analysis}
Let \(n\) be the number of requests currently stored inside the heap, and
\(k\) the parameter to \texttt{topKRequests()}.

\begin{table}[H]
\centering
\begin{tabular}{|p{4.3cm}|p{3.3cm}|p{2.5cm}|p{5.2cm}|}
\hline
\textbf{Operation} & \textbf{Time Complexity} & \textbf{Space Complexity} & \textbf{Explanation} \\
\hline
\texttt{swapNodes()}
& \(O(1)\) & \(O(1)\) &
Swaps two heap entries through a temporary. \\
\hline
\texttt{siftUp()},
\texttt{siftDown()}
& \(O(\log n)\) & \(O(1)\) &
Each one walks at most one root-to-leaf path. \\
\hline
\texttt{findIndexByID()}
& \(O(n)\) & \(O(1)\) &
Linear scan over the heap array. \\
\hline
\texttt{validateInput()}
& \(O(n^{2})\) & \(O(1)\) &
The duplicate-ID double loop is the dominant cost. \\
\hline
\texttt{buildStructure()}
& \(O(n)\) & \(O(n)\) &
Bottom-up heapify (Floyd's algorithm); the \(O(n)\) space holds the heap itself. \\
\hline
\texttt{dispatchNext()}
& \(O(\log n)\) & \(O(1)\) &
Replace the root with the last element and sift down. \\
\hline
\texttt{updatePriority()}
& \(O(n)\) & \(O(1)\) &
Dominated by the linear ID scan; the sift itself is \(O(\log n)\). \\
\hline
\texttt{sortRecords()}
& \(O(n \log n)\) average,
  \(O(n^{2})\) worst
& \(O(\log n)\) auxiliary &
Randomised Lomuto quicksort, fully in-place; the only extra memory is the recursion stack. \\
\hline
\texttt{topKRequests(k)}
& \(O(n + k \log n)\) & \(O(n)\) &
The live heap is copied into \texttt{tempHeap} so that the \(k\) extract-max operations do not modify it. \\
\hline
\end{tabular}
\caption{Task 2 time and space complexity.}
\end{table}

The auxiliary space of \texttt{sortRecords()} is only the recursion stack,
and no auxiliary array of size \(n\) is allocated --- this satisfies the
problem statement's in-place requirement.

\subsection{Test Case Demonstrations}

\begin{figure}[H]
\centering
\begin{subfigure}[b]{0.48\textwidth}
    \includegraphics[width=\textwidth,keepaspectratio]{task2_tc1.png}
    \caption{Task 2 -- Test Case 1 Output}
\end{subfigure}
\hfill
\begin{subfigure}[b]{0.48\textwidth}
    \includegraphics[width=\textwidth,keepaspectratio]{task2_tc2.png}
    \caption{Task 2 -- Test Case 2 Output}
\end{subfigure}

\vspace{0.6em}

\begin{subfigure}[b]{0.48\textwidth}
    \includegraphics[width=\textwidth,keepaspectratio]{task2_tc3.png}
    \caption{Task 2 -- Test Case 3 Output}
\end{subfigure}
\hfill
\begin{subfigure}[b]{0.48\textwidth}
    \includegraphics[width=\textwidth,keepaspectratio]{task2_tc4.png}
    \caption{Task 2 -- Test Case 4 Output}
\end{subfigure}
\caption{Task 2 test cases}
\end{figure}

\vspace{1em}

\noindent The test cases used in \texttt{main()} are hardcoded so that the
program output is deterministic; the screenshots above are produced by
compiling and running \texttt{TASK1.cpp} and \texttt{TASK2.cpp} directly.

% =====================================================================
\section{GitHub Repository}
% =====================================================================

The complete source code and the test-case screenshots for both tasks are
available in the following repository:
\begin{itemize}
    \item \url{https://github.com/dakshrathi-india/DSA-ASSIGNMENT-4}
\end{itemize}

\end{document}
